\section{Scientific motivation}
\label{sec:motivation}

The classical stellar--halo mass relation compresses a galaxy and its halo to
one number each,
\begin{equation}
  P(M_\star\mid M_{\rm h},z).
\end{equation}
That relation is foundational, but it omits two forms of structure that are
especially relevant for massive central galaxies. First, a halo has an assembly
history rather than only a present mass. Second, a galaxy has a radial stellar
mass distribution rather than only a total mass. The motivating object in
HongShao is therefore a profile-level, assembly-resolved extension,
\begin{equation}
  P\!\left[M_\star(<R,z)\given \Mpeak(t),\,\ctwo(t),\,z\right].
  \label{eq:ultimate_shmr}
\end{equation}
This is a conditional distribution: halos with identical measured inputs need
not host identical galaxies. The residual distribution is part of the result,
not a nuisance to be driven to zero.

This framing is motivated by the two-stage growth of massive galaxies. Early
dissipative growth preferentially builds dense central stellar structure,
whereas later mergers and stripping can add extended ex-situ envelopes. Halo
growth is not expected to encode every baryonic event, orientation, or merger
orbit, but it may encode the statistically predictable component of this
evolution. The broader galaxy--halo literature supports treating both mass and
assembly as possible predictors rather than assuming a scalar relation is
sufficient \citep{Wechsler2018}.

Earlier HongShao experiments established a clear statistical connection by
predicting principal-component coefficients of the curve of growth. PCA is a
powerful compression \citep{Jolliffe2002}, but an individual component is a
sample-defined weighted combination of radii. It does not immediately answer a
question such as ``does an earlier halo transition predict a larger central
mass fraction?'' The deposition-and-migration model attacks this interpretive
problem from the opposite direction, explicitly accumulating and moving mass,
but necessarily introduces assumptions about deposition kernels, efficiencies,
and radial migration.

Experiments 50 and 51 study a third route. They retain a statistical map but
replace opaque galaxy-side PCA coefficients with directly readable profile
coordinates. The central question is
\begin{equation}
  P\!\left(\vect{\theta}_{\cog}\given
  \vect{\theta}_{\rm MAH},\ctwo,z\right),
  \label{eq:central_question}
\end{equation}
where every element of $\vect{\theta}_{\cog}$ is a stellar mass, mass fraction,
or enclosed-mass radius. ``Physical'' is used here in a phenomenological sense:
the intermediate variables have physical meanings, but the fitted conversion
is not derived from conservation laws and is not claimed to be causal.

\subsection{Questions addressed by the two experiments}

\ExpFifty{} asks whether this readable conversion can match the established PCA
reference at one epoch, $z=0.4$. It separately tests whether halo assembly shape
and halo concentration carry information beyond final halo mass, whether the
model must be generative, and whether symbolic regression can expose a compact
nonlinear relation.

\ExpFiftyOne{} asks whether the result survives from $z=0.4$ to $z=2$. It also
separates two uses that are scientifically distinct: reconstructing the
progenitor history of a population selected at $z=0.4$, and painting galaxies
onto a halo catalog selected at the prediction epoch. The latter must not use
halo growth that has not yet occurred. Finally, \ExpFiftyOne{} tests whether a
simple temporal residual model yields coherent stochastic galaxy histories.

Throughout this report, every numerical comparison states the measured
quantity, its reference, and whether lower or higher values are preferable.
The intent is to make the logical chain auditable rather than to present only a
scoreboard.

